\section{Introduction}
\label{sec:intro}

Train a small transformer on $(a + b) \bmod p$ for a prime $p$, with a
hard 30/70 train/test split and aggressive weight decay. The training loss
collapses within a few hundred gradient steps, so the network has memorised
the training set. The test loss does not move. After tens of thousands of
\emph{additional} steps with no apparent progress, the test loss abruptly
collapses and the model generalises. \citet{power2022grokking} named this
phenomenon \emph{grokking} and showed that it occurs across a family of
small algorithmic tasks. \citet{nanda2023progress} subsequently showed
that during the apparent plateau the network is silently constructing a
discrete-Fourier-transform algorithm for the underlying group operation,
and that weight-decay-driven implicit regularisation is what tips the
optimiser from memorising circuits to generalising circuits.

Grokking sits at the intersection of three research frontiers:
generalisation theory, mechanistic interpretability, and optimisation
dynamics. The single-task setting is now reasonably well understood. The
\emph{multi-task} setting is not.

\paragraph{This paper.} We ask three questions:
\begin{enumerate}
\item When a single one-layer transformer is trained jointly on
      $\{+, -, \times\} \bmod p$, does grokking still occur on each task?
\item Does multi-task training accelerate, delay, or qualitatively change
      the per-task grokking dynamics relative to single-task training?
\item Are the discovered Fourier features shared across tasks or
      task-specific? In particular, do addition and subtraction (which
      both live on the additive group $\Z/p\Z$) develop a shared feature
      set, and does multiplication (which lives on the multiplicative
      group $(\Z/p\Z)^*$) develop a separate basis?
\end{enumerate}

\paragraph{Why these questions.}
All three operations have closed-form Fourier solutions (Section
\ref{sec:background}), so we can make falsifiable predictions about which
features should emerge. The multiplicative group $(\Z/p\Z)^*$ is cyclic of
order $p - 1$, which is a different cyclic group from the additive
$\Z/p\Z$ of order $p$, so the two groups have non-overlapping Fourier
bases. A single residual stream that solves all three tasks must therefore
either reuse one basis for both groups or carve out task-specific
sub-spaces. Both possibilities are interesting: the first would suggest
the model has discovered a representation that abstracts over group
structure, the second would suggest the model is internally
``time-multiplexing'' between two algorithms.

\paragraph{Contributions.}
\begin{itemize}
\item A clean replication of the \citet{nanda2023progress} single-task
      addition baseline at the published hyperparameters
      ($p = 113$, $\mathrm{train\_frac} = 0.30$, AdamW with
      $\mathrm{weight\_decay} = 1.0$).
\item A characterisation of per-task grokking dynamics in joint training
      on add+subtract and on add+subtract+multiply, with three random
      seeds for robustness.
\item A Fourier-feature analysis of the multi-task model's embedding
      matrix, comparing additive-group and multiplicative-group
      decompositions and quantifying overlap with the corresponding
      single-task models via cosine similarity.
\item Per-head ablations showing that attention heads partially
      specialise toward operation-conditional features.
\item A fully reproducible code release with config files, per-step
      metric histories, and all figures.\footnote{%
      \url{https://github.com/zaidanmir/grokking-multitask}}
\end{itemize}

The rest of the paper is structured as follows. Section
\ref{sec:background} reviews the algebraic and Fourier prerequisites and
the published single-task grokking story. Section \ref{sec:methods}
describes the model architecture, optimiser, and dataset construction.
Section \ref{sec:replication} reports the single-task replication.
Section \ref{sec:multitask} reports the multi-task experiments. Section
\ref{sec:analysis} reports the mechanistic Fourier analysis and head
ablations. Section \ref{sec:discussion} discusses what the findings imply
about transformer representations, and Section \ref{sec:future_work}
sketches open follow-ups.
